\chapter{Energy-parametrized \texorpdfstring{$L^2$}{L2} Basic}
\label{chap:canonical-basic}

The reusable Basic has three layers: intrinsic $L^2$ variance, extended energies, and a pure
optimization frontend.  It is parametrized by an energy and is not itself the choice of a
canonical energy.  None of these definitions assumes that the energy comes from a smooth gradient
or from a particular Dirichlet form.

\section{Variance and extended energy}

\begin{definition}[Intrinsic $L^2$ mean, centering, and variance]
  \label{def:l2-mean-variance}
  \lean{FunctionalInequalities.L2.mean,
        FunctionalInequalities.L2.center,
        FunctionalInequalities.L2.variance}
  \leanok
  For a real $L^2(\mu)$ class $f$, define its mean by the integral of any representative, its
  center by subtracting that mean, and
  \[
    \Var_\mu(f)=\norm{f-\textstyle\int f\dd\mu}_{L^2(\mu)}^2.
  \]
  All three objects are intrinsic to the almost-everywhere class.
\end{definition}

\begin{definition}[Extended energy]
  \label{def:extended-energy}
  \lean{FunctionalInequalities.Energy.Functional,
        FunctionalInequalities.Energy.Functional.FiniteAt,
        FunctionalInequalities.Energy.Functional.domain,
        FunctionalInequalities.Energy.Functional.finiteValue}
  \leanok
  An energy functional on a space $H$ is a map $\mathcal E:H\to[0,\infty]$.  Its domain is
  $\dom(\mathcal E)=\{f:\mathcal E(f)<\infty\}$; on this domain
  $\mathcal E_{\mathrm{fin}}$ is the corresponding finite nonnegative real value.
\end{definition}

\begin{definition}[Pre-energy and sequential relaxation]
  \label{def:preenergy-relaxation}
  \lean{FunctionalInequalities.Energy.PreEnergy,
        FunctionalInequalities.Energy.PreEnergy.DenseRealization,
        FunctionalInequalities.Energy.PreEnergy.ApproximatingSequence,
        FunctionalInequalities.Energy.PreEnergy.sequentialRelaxation}
  \leanok
  \uses{def:extended-energy}
  A pre-energy consists of a test space, a realization into an ambient pseudo-metric space, and a
  finite nonnegative energy on tests.  Its sequential relaxation at $f$ is the infimum of
  $\liminf_n\mathcal E_0(u_n)$ over realized test sequences converging to $f$; it is
  $\topenergy$ when no such sequence exists.
\end{definition}

\begin{theorem}[Quadratic relaxation is closed quadratic]
  \label{thm:quadratic-relaxation-closed}
  \lean{FunctionalInequalities.Energy.PreEnergy.Quadratic,
        FunctionalInequalities.Energy.ClosedQuadratic,
        FunctionalInequalities.Energy.PreEnergy.Quadratic.closedQuadratic_sequentialRelaxation}
  \leanok
  \uses{def:preenergy-relaxation}
  If a pre-energy has a linear realization, satisfies the quadratic laws on tests, and has dense
  realization, then its sequential relaxation is a densely defined lower-semicontinuous extended
  quadratic form satisfying the extended parallelogram identity.
\end{theorem}
\begin{proof}
  \leanok
  Recovery sequences supply finite approximants; linear combinations of approximating sequences
  transfer homogeneity and the parallelogram law through the liminf relaxation.  Lower
  semicontinuity follows from the diagonal characterization of sequential relaxation.
\end{proof}

\section{The optimization frontend}

\begin{definition}[Finite-domain test family]
  \label{def:poincare-test-family}
  \lean{FunctionalInequalities.Poincare.TestFamily,
        FunctionalInequalities.Poincare.TestFamily.rayleigh,
        FunctionalInequalities.Poincare.TestFamily.constant,
        FunctionalInequalities.Poincare.TestFamily.Holds}
  \leanok
  A test family provides a variance and a nonnegative energy on a chosen domain.  Its Rayleigh
  quotient uses the extended convention $0/0=0$ and $a/0=\topenergy$ for $a>0$.
  The optimal constant is the supremum of these quotients, and
  $\operatorname{Holds}(C)$ means that every test satisfies $\Var\le C\mathcal E$.
\end{definition}

\begin{definition}[Poincar\'e constant of an $L^2$ energy]
  \label{def:l2-energy-poincare-constant}
  \lean{FunctionalInequalities.Poincare.TestFamily.ofL2Energy,
        FunctionalInequalities.Poincare.L2Energy.constant,
        FunctionalInequalities.Poincare.L2Energy.Holds}
  \leanok
  \uses{def:l2-mean-variance,def:extended-energy,def:poincare-test-family}
  For a probability law $\mu$ and an extended energy $\mathcal E$ on real $L^2(\mu)$, let
  \[
    \CE(\mu,\mathcal E)
      =\sup_{f\in\dom(\mathcal E)}
        \frac{\Var_\mu(f)}{\mathcal E_{\mathrm{fin}}(f)}.
  \]
  The predicate $\operatorname{Poincare.L2Energy.Holds}(\mu,\mathcal E,C)$ is the corresponding inequality
  on every finite-energy $L^2$ class.
\end{definition}

\begin{theorem}[Energy-parametrized optimal constant order law]
  \label{thm:l2-energy-holds-order}
  \lean{FunctionalInequalities.Poincare.L2Energy.holds_iff,
        FunctionalInequalities.Poincare.L2Energy.holds_iff_constant_le}
  \leanok
  \uses{def:l2-energy-poincare-constant}
  For every finite nonnegative constant $C$,
  \[
    \operatorname{Poincare.L2Energy.Holds}(\mu,\mathcal E,C)
    \quad\Longleftrightarrow\quad
    \CE(\mu,\mathcal E)\le C.
  \]
\end{theorem}
\begin{proof}
  \leanok
  This is the order characterization of the supremum of the finite-domain Rayleigh quotients,
  including zero energy and infinite optimal constants.
\end{proof}


\begin{lemma}[Relaxation preserves Poincar\'e constants]
  \label{lem:preenergy-poincare-relaxation}
  \lean{FunctionalInequalities.Poincare.TestFamily.ofPreEnergy,
        FunctionalInequalities.Poincare.TestFamily.holds_ofPreEnergy_iff,
        FunctionalInequalities.Poincare.TestFamily.constant_ofPreEnergy_eq}
  \leanok
  \uses{def:l2-mean-variance,def:preenergy-relaxation,
        def:poincare-test-family,def:l2-energy-poincare-constant}
  Let $\mu$ be a probability measure on any measurable space and $P$ a finite pre-energy
  realized in real $L^2(\mu)$. Give each concrete test its realized variance and raw energy.
  For every finite $C\ge0$, its Poincar\'e inequality holds at $C$ exactly when the inequality
  holds on the finite-energy domain of $P.\mathrm{sequentialRelaxation}$ at $C$.
  The two extended optimal constants agree. No density or quadraticity is required; empty test
  classes, zero constants and infinite optima are included.
\end{lemma}
\begin{proof}
  \leanok
  Constant sequences give one direction. At every finite-energy point, a recovery sequence
  and continuity of $L^2$ variance give the converse, also for $C=0$. The existing finite-upper-bound
  characterization in \texttt{ENNReal} identifies the extended optima.
\end{proof}
